\setcounter{secnumdepth}{3}
\titlespacing*{\chapter}{0pt}{-0.1cm}{1cm}
\titleformat{\chapter}[display]
  {\bfseries\LARGE}  
  {Phần~\thechapter}{0pt}{\raggedright}

\chapter{Các thuật toán suy diễn}
\section{Forward Chaining}

\subsection{Cơ chế hoạt động}

Thuật toán được cài đặt theo hướng sử dụng \textbf{hàng đợi (agenda)} để quản lý các fact cần xử lý. Đồng thời, nhóm sử dụng thêm một cấu trúc chỉ mục (tương tự hash map) để nhanh chóng xác định các luật bị ảnh hưởng khi một fact mới xuất hiện.

Ý tưởng chính của Forward Chaining là lặp lại quá trình suy diễn cho đến khi đạt trạng thái ổn định (không sinh thêm fact mới). Cụ thể:

\begin{itemize}
    \item Khởi tạo agenda với các fact ban đầu.
    \item Lặp lại:
    \begin{itemize}
        \item Lấy một fact từ agenda.
        \item Cập nhật các luật liên quan (giảm \texttt{count}).
        \item Nếu một luật thỏa mãn (\texttt{count} = 0), suy ra fact mới.
        \item Thêm fact mới vào KB và đưa vào agenda.
    \end{itemize}
    \item Dừng khi agenda rỗng (đạt điểm cố định).
\end{itemize}

Sau bước này:
\begin{itemize}
    \item Nếu tất cả các ô đã được xác định giá trị $\rightarrow$ kết thúc.
    \item Ngược lại $\rightarrow$ chuyển sang Backtracking để tiếp tục tìm nghiệm.
\end{itemize}

\subsection{Kết hợp Backtracking và heuristic MRV}

Khi Forward Chaining không thể suy diễn thêm (đạt trạng thái ổn định) nhưng bài toán vẫn chưa giải xong, nhóm sử dụng \textbf{Backtracking} để tìm kiếm phần còn lại của nghiệm.

Để giảm không gian tìm kiếm, thuật toán sử dụng heuristic \textbf{MRV (Minimum Remaining Values)}. Thay vì chọn ô theo thứ tự cố định, MRV ưu tiên chọn ô có domain nhỏ nhất (tức là ít giá trị khả thi nhất).

Ý tưởng của MRV:
\begin{itemize}
    \item Ô có domain nhỏ thường bị ràng buộc nhiều hơn.
    \item Dễ phát hiện mâu thuẫn sớm nếu chọn sai.
    \item Giảm số nhánh cần duyệt trong cây tìm kiếm.
\end{itemize}

Quy trình Backtracking kết hợp MRV như sau:

\begin{itemize}
    \item Chọn ô $(i,j)$ chưa xác định với domain nhỏ nhất.
    \item Thử lần lượt từng giá trị $v$ trong domain$(i,j)$:
    \begin{itemize}
        \item Gán $(i,j) = v$ và thêm fact \texttt{Val(i,j,v)} vào KB.
        \item Chạy lại Forward Chaining để lan truyền ràng buộc.
        \item Nếu xảy ra mâu thuẫn (domain rỗng hoặc conflict) thì quay lui và thử giá trị khác.
    \end{itemize}
\end{itemize}

Nhờ kết hợp Forward Chaining và MRV, thuật toán vừa tận dụng được khả năng suy diễn để thu hẹp domain, vừa đảm bảo tìm được nghiệm đầy đủ với số bước thử ít hơn so với Backtracking thuần túy.

\subsection{Ví dụ minh họa}

Xét bài toán Futoshiki kích thước $3 \times 3$ với dữ liệu đầu vào như sau (intput1):

\textbf{Grid:}
\[
\begin{bmatrix}
0 & 2 & 0 \\
0 & 0 & 0 \\
0 & 0 & 0
\end{bmatrix}
\]

\textbf{Ràng buộc ngang:}
\begin{itemize}
    \item $(1,1) < (1,2)$
\end{itemize}

\textbf{Ràng buộc dọc:}
\begin{itemize}
    \item $(1,2) > (2,2)$
    \item $(2,1) > (3,1)$
\end{itemize}

\textbf{Bước 1: Khởi tạo KB và Domain}

\begin{itemize}
    \item Từ dữ kiện ban đầu: \texttt{Given(1,2,2)} $\Rightarrow$ \texttt{Val(1,2,2)}
    \item Domain ban đầu:
    \[
    (1,2) = \{2\}, \quad \text{các ô còn lại} = \{1,2,3\}
    \]
\end{itemize}

\textbf{Bước 2: Forward Chaining}

Từ \texttt{Val(1,2,2)}, suy ra:

\begin{itemize}
    \item Loại bỏ $2$ khỏi hàng 1: $(1,1),(1,3)$ không thể là $2$
    \item Loại bỏ $2$ khỏi cột 2: $(2,2),(3,2)$ không thể là $2$
\end{itemize}

Xét ràng buộc $(1,1) < (1,2)$ với $(1,2) = 2$:
\begin{itemize}
    \item $(1,1) < 2 \Rightarrow (1,1) = 1$
\end{itemize}

Suy ra:
\[
\texttt{Val(1,1,1)}
\]

Tiếp tục Forward Chaining:
\begin{itemize}
    \item Loại $1$ khỏi hàng 1 và cột 1
\end{itemize}
Suy ra:
\[
\texttt{Val(1,3,3)}
\]

Sau một số bước suy diễn, domain được thu hẹp nhưng vẫn tồn tại trường hợp chưa xác định hết các ô.

\textbf{Bước 3: Áp dụng MRV}

Xét các ô chưa xác định, chọn ô có domain nhỏ nhất. Ví dụ:
\begin{itemize}
    \item Giả sử ô $(2,2)$ còn domain $\{1,3\}$ (nhỏ nhất)
\end{itemize}

Chọn ô $(2,2)$ để thử giá trị.

\textbf{Bước 4: Backtracking + Forward Chaining}

\begin{itemize}
    \item Thử $(2,2) = 1$:
    \begin{itemize}
        \item Thêm \texttt{Val(2,2,1)} vào KB
        \item Chạy lại Forward Chaining
        \item Nếu xảy ra mâu thuẫn $\rightarrow$ quay lui
    \end{itemize}
    
    \item Nếu không hợp lệ, thử $(2,2) = 3$:
    \begin{itemize}
        \item Tiếp tục suy diễn và lan truyền ràng buộc
    \end{itemize}
\end{itemize}

Quá trình này tiếp tục cho đến khi tất cả các ô đều được gán giá trị.

\subsection{Mã giả thuật toán}
\begin{algorithm}[H]
\footnotesize
\caption{FC Solver cho Futoshiki}
\KwIn{Dữ liệu bài toán $data$}
\KwOut{Nghiệm $solution$}

$domain \leftarrow$ domain\_init($data$) \;
$kb \leftarrow$ build\_kb($data$) \;

$result \leftarrow$ backtrack($kb$, $domain$) \;

\If{$result = \varnothing$}{
    \Return $\varnothing$ \;
}

\For{mỗi ô $(i,j)$}{
    $v \leftarrow$ phần tử duy nhất trong $result(i,j)$ \;
    $data.grid[i-1][j-1] \leftarrow v$ \;
}

\Return $data$ \;
\end{algorithm}

\begin{algorithm}[H]
\footnotesize
\caption{Backtracking với FC + MRV}
\KwIn{$kb$, $domain$}
\KwOut{Domain hoàn chỉnh hoặc $\varnothing$}

$success \leftarrow$ propagate($kb$, $domain$) \;

\If{NOT $success$}{
    \Return $\varnothing$ \;
}

Kiểm tra tất cả ô đã xác định chưa \;

\If{mọi ô chỉ có 1 giá trị}{
    \Return $domain$ \;
}

$cell \leftarrow$ select\_mrv($domain$) \;

\If{$cell = \varnothing$}{
    \Return $\varnothing$ \;
}

$(i,j) \leftarrow cell$ \;
$values \leftarrow$ các giá trị trong $domain(i,j)$ (đã sắp xếp) \;

\For{mỗi $v \in values$}{

    $new\_kb \leftarrow$ copy($kb$) \;
    $new\_domain \leftarrow$ deepcopy($domain$) \;

    $new\_domain(i,j) \leftarrow \{v\}$ \;
    Thêm $Val(i,j,v)$ vào $new\_kb$ \;

    $result \leftarrow$ backtrack($new\_kb$, $new\_domain$) \;

    \If{$result \neq \varnothing$}{
        \Return $result$ \;
    }
}

\Return $\varnothing$ \;
\end{algorithm}

\begin{algorithm}[H]
\footnotesize
\caption{Chọn biến theo MRV}
\KwIn{$domain$}
\KwOut{Ô $(i,j)$ có domain nhỏ nhất}

$best\_cell \leftarrow \varnothing$ \;
$best\_size \leftarrow \infty$ \;

\For{mỗi ô $cell$}{
    $values \leftarrow domain(cell)$ \;

    \If{$|values| > 1$ và $|values| < best\_size$}{
        $best\_size \leftarrow |values|$ \;
        $best\_cell \leftarrow cell$ \;
    }
}

\Return $best\_cell$ \;
\end{algorithm}

\begin{algorithm}[H]
\footnotesize
\caption{Propagation (Forward Checking)}
\KwIn{$kb$, $domain$}
\KwOut{True nếu thành công, False nếu mâu thuẫn}

$step \leftarrow 0$ \;

\While{True}{

    $old\_domain \leftarrow$ deepcopy($domain$) \;

    forward\_chaining($kb$) \;

    \If{phát hiện xung đột}{
        \Return False \;
    }

    Cập nhật $domain$ từ $kb$ \;

    \If{có domain rỗng}{
        \Return False \;
    }

    \If{$domain = old\_domain$ và không có fact Val mới}{
        \textbf{break}
    }
}

\Return True \;
\end{algorithm}

\subsection{Nhận xét}

Forward Chaining giúp lan truyền ràng buộc rất nhanh và giảm đáng kể không gian tìm kiếm ban đầu. Tuy nhiên, do chỉ dựa vào suy diễn nên thuật toán có thể bị dừng sớm khi không còn luật nào áp dụng được.

Vì vậy, việc kết hợp với Backtracking là cần thiết để đảm bảo tìm được nghiệm đầy đủ cho mọi test case.


\section{Backward Chaining}

Nếu Forward Chaining là phương pháp suy diễn tiến từ các dữ kiện ban đầu để sinh ra các sự thật mới, thì \textbf{Backward Chaining} là cách suy diễn theo hướng ngược lại: bắt đầu từ một mục tiêu cần kiểm tra và truy ngược về các luật và sự thật trong cơ sở tri thức (KB).

Trong đồ án này, Backward Chaining không được dùng để giải toàn bộ bảng, mà chủ yếu đóng vai trò \textbf{hỗ trợ truy vấn và kiểm tra tính khả thi} của một giá trị tại một ô cụ thể.

\subsection{Cơ chế hoạt động}

Thuật toán được cài đặt dưới dạng hàm đệ quy \texttt{backward\_chain(kb, goal, visited, memo)}. Với một mục tiêu cần chứng minh, thuật toán thực hiện:

\begin{itemize}
    \item Nếu mục tiêu đã tồn tại trong tập \texttt{facts} $\rightarrow$ trả về \texttt{True}.
    
    \item Ngược lại, tìm các luật có kết luận trùng với mục tiêu thông qua cấu trúc \textbf{conclusion index} (giúp truy xuất nhanh thay vì duyệt toàn bộ luật).
    
    \item Với mỗi luật tìm được, lần lượt chứng minh các tiền đề của luật bằng cách gọi đệ quy.
    
    \item Nếu tất cả tiền đề đều đúng $\rightarrow$ mục tiêu đúng.
\end{itemize}

Để tối ưu, thuật toán sử dụng thêm hai kỹ thuật:

\begin{itemize}
    \item \textbf{Loop detection (\texttt{visited}):} Tránh lặp vô hạn khi các luật tạo thành vòng.
    
    \item \textbf{Memoization (\texttt{memo}):} Lưu kết quả các mục tiêu đã chứng minh để tránh tính lại nhiều lần.
\end{itemize}

\subsection{Ứng dụng trong truy vấn giá trị}

Trong bài toán Futoshiki, Backward Chaining được dùng để xác định các giá trị khả thi cho một ô $(i,j)$.

Thay vì chứng minh trực tiếp $Val(i,j,v)$, nhóm sử dụng cách tiếp cận đơn giản hơn:

\begin{itemize}
    \item Với mỗi giá trị $v \in \{1,\dots,N\}$:
    \begin{itemize}
        \item Kiểm tra xem có chứng minh được $NotVal(i,j,v)$ hay không.
    \end{itemize}
    
    \item Nếu chứng minh được $\rightarrow$ loại $v$ khỏi domain.
    
    \item Nếu không chứng minh được $\rightarrow$ giữ $v$ là giá trị khả thi.
\end{itemize}

Cách làm này giúp loại bỏ các giá trị không hợp lệ mà không cần thử toàn bộ không gian tìm kiếm.

\subsection{Ví dụ minh họa}

Xét bài toán Futoshiki kích thước $3 \times 3$ với dữ liệu:

\textbf{Grid:}
\[
\begin{bmatrix}
0 & 2 & 0 \\
0 & 0 & 0 \\
0 & 0 & 0
\end{bmatrix}
\]

\textbf{Ràng buộc:}
\begin{itemize}
    \item $(1,1) < (1,2)$
    \item $(1,2) > (2,2)$
    \item $(2,1) > (3,1)$
\end{itemize}

Ta cần kiểm tra xem ô $(1,1)$ có thể nhận giá trị $2$ hay không.

\textbf{Bước 1: Thiết lập truy vấn}

Thay vì chứng minh trực tiếp $Val(1,1,2)$, ta kiểm tra:
\[
\texttt{backward\_chain}(kb, NotVal(1,1,2))
\]

\textbf{Bước 2: Tìm luật liên quan}

Từ KB, tồn tại luật ràng buộc hàng:
\[
Val(1,2,2) \Rightarrow NotVal(1,1,2)
\]

\textbf{Bước 3: Chứng minh mục tiêu con}

Để chứng minh $NotVal(1,1,2)$, ta cần chứng minh:
\[
Val(1,2,2)
\]

Từ dữ kiện đề bài:
\[
Given(1,2,2) \Rightarrow Val(1,2,2)
\]

Vì $Given(1,2,2)$ là fact $\rightarrow$ suy ra $Val(1,2,2)$ đúng.

\textbf{Bước 4: Kết luận}

Suy ra:
\[
NotVal(1,1,2) \text{ là True}
\]

Do đó, giá trị $2$ bị loại khỏi domain của ô $(1,1)$.

\textbf{Kết luận cuối cùng:} Ô $(1,1)$ \textbf{không thể} nhận giá trị $2$.

\textbf{Nhận xét:}

Ví dụ trên cho thấy Backward Chaining có thể nhanh chóng loại bỏ các giá trị không hợp lệ bằng cách truy ngược từ mục tiêu. Thay vì thử gán trực tiếp, thuật toán chỉ cần kiểm tra các luật liên quan trong KB, giúp giảm đáng kể số phép thử trong quá trình giải bài toán.

\textbf{Kết luận:} Ô $(1,2)$ không thể nhận giá trị $1$.

\subsection{Mã giả thuật toán}

\subsection{Nhận xét}

Backward Chaining phù hợp cho các truy vấn cục bộ, đặc biệt khi cần kiểm tra nhanh tính hợp lệ của một giá trị. 

Trong khi đó, Forward Chaining hiệu quả hơn trong việc lan truyền ràng buộc trên toàn bộ bảng. Việc kết hợp hai phương pháp giúp hệ thống vừa giảm được không gian tìm kiếm, vừa đảm bảo tính chính xác khi suy diễn.

\section{Thuật toán A*}


A* là một thuật toán tìm kiếm có thông tin (informed search), kết hợp giữa:
\begin{itemize}
    \item Chi phí đã đi từ trạng thái ban đầu đến trạng thái hiện tại.
    \item Một hàm ước lượng chi phí còn lại để đạt tới trạng thái đích.
\end{itemize}

Tại mỗi bước, A* sẽ chọn trạng thái có giá trị nhỏ nhất theo hàm:
\[
f(n) = g(n) + h(n)
\]
trong đó:
\begin{itemize}
    \item $g(n)$: chi phí từ trạng thái ban đầu đến trạng thái $n$.
    \item $h(n)$: hàm heuristic ước lượng chi phí từ $n$ đến trạng thái đích.
\end{itemize}

Trong bối cảnh bài toán Futoshiki, mỗi trạng thái tương ứng với một cấu hình của bảng (một phần hoặc toàn bộ các ô đã được gán giá trị). Quá trình tìm kiếm được thực hiện bằng cách:
\begin{itemize}
    \item Chọn một ô chưa được gán giá trị.
    \item Thử gán một giá trị hợp lệ từ miền giá trị của ô đó.
    \item Sinh ra trạng thái mới và đưa vào hàng đợi ưu tiên.
\end{itemize}

Thuật toán sử dụng một \textbf{priority queue} (hàng đợi ưu tiên) để luôn mở rộng trạng thái có giá trị $f(n)$ nhỏ nhất trước. Đồng thời, một tập \textbf{visited} được sử dụng để tránh việc lặp lại các trạng thái đã xét.

A* đảm bảo tìm được nghiệm tối ưu nếu hàm heuristic thỏa mãn tính chất admissible. Tuy nhiên, hiệu quả của thuật toán phụ thuộc rất lớn vào chất lượng của heuristic được sử dụng. 